\section{Introduction}
\IEEEPARstart{L}{arge} language model (LLM) agents increasingly participate in
industrial operation: they generate operator-reviewed OPC UA control sequences
\cite{huang2026made}, invoke factory skills over tool protocols
\cite{koehle2026llmdriven}, coordinate shop-floor tasks
\cite{lim2025llmenabled,zhao2026shopfloor}, and schedule a simulated
robot-operated line \cite{waseem2026pretrained}. Their usefulness depends on
connecting task reasoning to the equipment and production context in which
a proposal will be executed. That connection must express who may change a
parameter, which product and recipe are active, what was actually commanded,
and how the process responded.

Consider an injection-molding task that asks an agent to bring part weight
into specification. The agent may consult recent measurements, diagnose
insufficient packing, and propose higher holding pressure. Executing this
proposal also requires a writable signal mapping, applicable parameter
limits, an authorized production scope, and the required approval. The
subsequent record must distinguish a pressure request, its endpoint
confirmation, and the measured weight. These are different events with
different evidence, even when one task spans the complete operation.

This information spans agent identity, industrial I/O, production
management, and process history. Three distinctions must survive this
handoff: permission to invoke a tool versus authority over a production
object, transport acceptance versus endpoint confirmation, and a trial
command versus a retained recipe. Reconstructing these distinctions inside
every agent application couples execution policy to the selected reasoning
engine and obscures how an intervention relates to a running batch. We
address the following systems question: how can an industrial
agent platform retain a common execution and evidence contract when the
proposal strategy, equipment protocol, and production context change?

We present a governed supervisory execution framework connecting agent
teams to semantic data-acquisition (DAQ) and write-control (DCW) nodes.
A deterministic execution boundary sits beneath the replaceable strategy
layer. The strategy acquires evidence, proposes a change, and interprets
observations. The platform resolves authority, applies source-specific
limits and required approval, dispatches commands, and grades confirmation.
Line--product--recipe--run objects supply production context.
Separate task, intervention, and batch records then connect the original
intent to commanded controls and observed outputs.

The contribution is an integration of production-scoped authority, graded
command confirmation, and route-specific observation records into a common
execution-and-evidence contract for industrial agent supervision:
\begin{enumerate}
\item a capability-based architecture that binds channel identity and
semantic industrial nodes to line--product--recipe--run context, separating
information exchange from execution authority;
\item route-specific execution mechanisms that join admission, approval,
graded confirmation, observation, and compensation while preserving
distinct node-intervention and recipe-batch records; and
\item an implementation and evaluation of this contract across replaceable
strategies and protocol mappings, using admission ablation, confirmation
drills, process missions, repeated regression runs, and a recipe trajectory.
\end{enumerate}

Sections~\ref{sec:arch}--\ref{sec:mech} explain the architecture and execution
mechanisms; Sections~\ref{sec:eval}--\ref{sec:limits} evaluate and interpret
their simulator-backed evidence.
